% Part: history 
% Chapter: biographies 
% Section: julia-robinson
\documentclass[../../../include/open-logic-section]{subfiles}

\begin{document}

\olfileid{his}{bio}{rob}

\olsection{Robinson}

\olphoto{robinson-julia}{Robinson}

鲁滨逊是美国数学家。她主要因在判定问题方面的工作而闻名，其中最著名的是她对解决 希尔伯特 第十问题所作的贡献。鲁滨逊 于 1919 年 12 月 8 日出生于密苏里州圣路易斯。鲁滨逊 曾回忆，她童年时就已经对数字着迷\citep[4]{Reid1986}。九岁时，她患上了猩红热，随后又多次复发风湿热。这迫使她长时间卧床，致使学业落后。尽管在家庭教师的帮助下她得以赶上进度，但疾病留下的身体后遗症影响了她的一生。

尽管童年饱受病痛折磨，鲁滨逊 高中毕业时仍在数学和科学方面获得了多项奖励。她在圣地亚哥州立学院开始大学学业，并于大四时转学至加州大学伯克利分校。在那里，她受到了数学家 鲁滨逊的影响。他们成为了好朋友，并于 1941 年结婚。作为教员的配偶，鲁滨逊 被禁止在伯克利数学系任教。尽管她继续旁听数学课程，但她原本希望离开大学去组建家庭。然而，婚礼后不久，鲁滨逊 患了肺炎。她被告知，由于童年时患过的风湿热，她的心脏上积累了大量疤痕组织。由于疤痕组织非常严重，医生预测她难以活过四十岁，并建议她不宜生育\citep[13]{Reid1986}。

鲁滨逊 陷入了长期的消沉，但最终决定继续研究数学。她回到伯克利，于 1948 年在 塔斯基的指导下完成博士学业并获博士学位。实数的一阶理论已被 塔斯基 证明是可判定的，而由 哥德尔的工作可知，自然数的一阶理论是不可判定的。有理数的一阶理论是否可判定，曾是一个重大的未解问题。在其博士论文中 \citeyearpar{Robinson1949}，鲁滨逊 证明了其不可判定。

出于对判定问题的兴趣，鲁滨逊 随后尝试寻找 希尔伯特第十问题的解法。该问题是 David 希尔伯特 于 1900 年提出的著名 23 个数学问题之一。第十问题问的是，是否存在一种算法，能够在有限时间内判定一个整系数多项式方程（例如 $3x^2 - 2y +3 = 0$）是否有整数解。此类问题被称为\emph{丢番图问题}。在取得一些初步成功后，鲁滨逊 与同样在研究该问题的 Martin Davis（马丁·戴维斯）和 Hilary Putnam（希拉里·普特南）联手。他们成功证明了指数丢番图问题（其中未知数也可能出现在指数位置）是不可判定的，并证明了某个猜想（后来被称为“J.R.”）蕴涵 希尔伯特 第十问题是不可判定的\citep{DavisPutnamRobinson1961}。鲁滨逊 在整个 1960 年代持续研究该问题。1970 年，年轻的俄罗斯数学家 Yuri Matijasevich（尤里·马季亚谢维奇）最终证明了 J.R. 猜想。上述综合结果现在被称为 Matijasevich--鲁滨逊--Davis--Putnam 定理，简称 MRDP 定理。Matijasevich 和 鲁滨逊 成为了朋友，并合作发表了多篇论文。在致 Matijasevich 的一封信中，鲁滨逊 曾写道：“事实上，我非常高兴，尽管我们相隔数千英里，但合作取得的进展显然比我们任何一人单独研究所取得的进展都要多”\citep[45]{Matijasevich1992}。

鲁滨逊 是美国数学学会的首位女性主席，也是首位当选美国国家科学院院士的女性。确诊白血病后，她于 1985 年 7 月 30 日去世，享年 65 岁。

\begin{reading}
鲁滨逊 的数学论文收录在其\emph{全集}\citep{Robinson1996}中，该书还重印了她的美国国家科学院传记回忆录\citep{Feferman1994}。鲁滨逊 的姐姐 Constance Reid（康斯坦丝·里德）根据访谈出版了一部《Julia 的自传》\citep{Reid1986}，以及一部详尽的回忆录\citep{Reid1996}。一部关于 鲁滨逊 和 希尔伯特 第十问题的纪录短片由 George Csicsery（乔治·齐克塞里）执导\citep{Csicsery2016}。关于 Yuri Matijasevich 与 鲁滨逊 的合作以及她对其工作的影响，可参阅简短回忆文章\citep{Matijasevich1992}。
\end{reading}

\end{document}